\section{The environment}\label{sec:env}

SharpeArena's Cargo workspace has two members. The crate \texttt{sharpearena} holds everything determinism-critical: the stepping engine, the wire contract, the scenario generator, the limit-order-book matching engine, mandates and the execution-noise model, built on the published SharpeBench engine crates rather than a vendored copy. The member \texttt{sharpearena-wasm} compiles the same kernel to WebAssembly and tests it against the native engine. A third Rust crate, \texttt{sharpearena-py}, is explicitly excluded from that workspace and built separately with maturin; it binds the engine into a Python package that adds the ecosystem adapters: Gymnasium \citep{towers2024gymnasium}, PettingZoo \citep{terry2021pettingzoo}, Minari \citep{minari}, a \texttt{verifiers} training environment, and an MCP server. The adapters are thin by policy; anything that must be byte-identical lives in Rust.

\subsection{Point-in-time observation, and the guard behind it}\label{sec:leakfree}

The engine owns the time cursor. On each step it emits a \texttt{MarketObservation} whose per-symbol snapshot carries a trailing window of closes up to the decision date, with optional fundamentals and news channels derived from that same trailing window. There is no API on the data layer that returns a bar after the cursor, so a policy cannot observe a future bar. Discipline inside generated code does not give the same guarantee: on the causal task category of AlphaQT-Bench, about half of the twelve evaluated models wrote some executable factor code that uses future data \citep[p.~7]{luo2026alphaqt}. The guarantee bounds what can be requested, not what can be inferred from the trailing window a request returns (see \cref{sec:limitations}). A property test iterates every scenario tier and asserts that no surfaced window extends past the cleared bar and that each window is a suffix of the true history.

One layer out, at the agent-harness boundary, sits \texttt{LookaheadGuard}, a refusal layer for tools and policies that try to read what the environment will not give. It carries a deny-list of named operations, each mapped to the reason it is refused: reading the raw dataset (the answer key), slicing past the current index, peeking the next bar, or feeding the scenario seed to the policy as a feature, which would identify the episode and trivialize generalization. Substring tells catch novel operation names, and an index check bounds any point-in-time slice: reading an index strictly greater than the current bar raises. A research escape hatch exists as an explicit environment variable, so relaxing the guard is a recorded decision rather than an accident. The guard is defense in depth behind the structural guarantee: a harness that bypasses it still cannot make the environment leak.

\subsection{Observation and action spaces}\label{sec:spaces}

An implementer sees two JSON documents and nothing else. \Cref{tab:obs-space} gives the observation, \cref{tab:act-space} the decision; both schemas are draft 2020-12, both close the object with \texttt{additionalProperties: false}, and both are committed beside the conformance kit of \cref{sec:contract}.

The observation is flat. \texttt{date} is the clock: an ISO-8601 date naming the decision point, and the bound every other field respects. \texttt{cash} is the agent's own cash. \texttt{symbols} carries one snapshot per tradable instrument, and \texttt{portfolio} the agent's own holdings, never a peer's pending state. A snapshot's \texttt{close\_history} is the trailing window: the last $L$ closes up to and including \texttt{date}, oldest first, where $L$ is the scenario's disclosure setting. The default is 20 closes with no optional channels; the shipped presets are a data-poor tier at 3 closes and a data-rich tier at 50 closes plus fundamentals and news. Both optional channels are derived from the surfaced window itself and therefore inherit its cutoff: \texttt{fundamentals} is a named map computed from those closes (trailing return, window high, window low) and \texttt{news} is one deterministic headline per symbol whose wording is a function of the window's trailing return. The v1.0 observation carries no mandate field. The mandate is a scenario-level object (\cref{sec:mandates}) rendered into the episode prompt by the Python task layer, so the fixed reference policies of \cref{sec:findings} act mandate-blind, which is what the failure taxonomy of \cref{sec:f7} counts.

\begin{table}[t]
\centering
\caption{The \texttt{MarketObservation} wire object at contract v1.0. Optional fields default to empty and are absent from the schema's \texttt{required} list by construction, which is what makes an additive field backwards-compatible.}
\label{tab:obs-space}
\footnotesize
\begin{tabular}{@{}llll@{}}
\toprule
Field & Type & Presence & Meaning \\
\midrule
\multicolumn{4}{@{}l}{\emph{Object} \texttt{MarketObservation}} \\
\texttt{date} & string (ISO-8601) & required & decision point; the point-in-time cutoff \\
\texttt{cash} & number & required & cash available to this agent \\
\texttt{symbols} & array of \texttt{SymbolSnapshot} & required & one snapshot per tradable instrument \\
\texttt{portfolio} & array of \texttt{PositionState} & required & this agent's holdings; empty at a cold start \\
\addlinespace
\multicolumn{4}{@{}l}{\emph{Object} \texttt{SymbolSnapshot}} \\
\texttt{symbol} & string & required & instrument identifier \\
\texttt{close\_history} & array of number & required & trailing closes through \texttt{date}, oldest first \\
\texttt{fundamentals} & map string to number & optional & derived fields; empty by default \\
\texttt{news} & array of string & optional & derived headlines; empty by default \\
\addlinespace
\multicolumn{4}{@{}l}{\emph{Object} \texttt{PositionState}} \\
\texttt{symbol} & string & required & instrument identifier \\
\texttt{shares} & number & required & signed holding; negative is a short \\
\texttt{avg\_price} & number & required & displayed average entry price \\
\bottomrule
\end{tabular}
\end{table}

The action space is a signed target-weight vector over the observed universe. A \texttt{Decision} is a list of per-instrument orders plus an optional free-text \texttt{reasoning} string that the trajectory records. An empty order list is a hold. Within an order, \texttt{action} is a label drawn from a closed four-variant enum (\texttt{buy}, \texttt{sell}, \texttt{hold}, \texttt{close}) and carries no sizing; all sizing is in \texttt{target\_weight}, which is a signed target portfolio weight for that symbol. The sign is the shorting semantics: a negative \texttt{target\_weight} is a short position, and the \texttt{signed-weight-short} conformance fixture exercises exactly that case, pairing \texttt{action: sell} with \texttt{target\_weight: -0.3} while a sibling order holds at $0.0$ and a third buys at $0.5$. Adding a fifth action variant is classified as breaking (\cref{sec:contract}) precisely because a consumer that matches the enum exhaustively would misparse it.

Four bounds are asserted by the conformance kit rather than by the schema's type system: every \texttt{action} lies in the enum, every \texttt{target\_weight} is finite, every order \texttt{symbol} is drawn from the universe the observation actually offered, and every \texttt{confidence} lies in $[0,1]$. The schema places no numeric bound on \texttt{target\_weight} itself; the gross-exposure ceiling an episode is graded against is a property of its sampled mandate (\cref{sec:mandates}). A legacy decision that omits \texttt{confidence}, \texttt{rationale} and \texttt{reasoning} still parses and is still well-formed, and the \texttt{legacy-decision-shape} fixture asserts it on every commit. Because both schemas close the object, a field the engine emits but the schema omits is an interoperability break rather than a documentation gap, so the conformance kit compares the serialized keys of fully-populated wire values against the schema's declared properties in both directions, for the observation and the decision alike.

A stated \texttt{confidence} feeds only a calibration diagnostic, and the Python paths count only confidences an agent actually stated. The local field runner records a confidence--outcome pair only for a decision with at least one order that states a confidence, using the mean over those orders, and pairs it with the sign of the next step's reward. The next step is the right one because the reward booked at a step is the price move on the holdings the previous decision chose, plus that step's trading cost; a run's final decision therefore adds no pair. \texttt{score\_run} scores a bare return track, which states no confidence, so it reports no calibration Brier score and zero calibration observations rather than scoring a filled-in $0.5$ on every bar. Rust paths that run the pinned SharpeBench 0.27.0 simulator directly, \texttt{run\_backtest\_checked} among them, keep that release's rule: an omitted confidence reads as $0.5$, as the wire default in \cref{tab:act-space} states, and every step, holds included, is paired with the same step's return. The two rules coexist until a SharpeBench release carries the new one and the pin here moves to it. No table or figure in this paper reports a calibration statistic.

\begin{table}[t]
\centering
\caption{The \texttt{Decision} wire object at contract v1.0. Optional fields carry the stated default, so a decision written against an earlier revision of the surface still parses.}
\label{tab:act-space}
\footnotesize
\begin{tabular}{@{}llll@{}}
\toprule
Field & Type & Presence & Meaning \\
\midrule
\multicolumn{4}{@{}l}{\emph{Object} \texttt{Decision}} \\
\texttt{orders} & array of \texttt{Order} & required & per-instrument instructions; empty is a hold \\
\texttt{reasoning} & string & optional (\texttt{""}) & rationale captured into the trajectory \\
\texttt{cost} & \texttt{DecisionCost} & optional (absent) & self-reported spend for this decision \\
\addlinespace
\multicolumn{4}{@{}l}{\emph{Object} \texttt{Order}} \\
\texttt{symbol} & string & required & identifier; must be an observed symbol \\
\texttt{action} & enum & required & \texttt{buy}, \texttt{sell}, \texttt{hold} or \texttt{close}; label only \\
\texttt{target\_weight} & number & required & signed target weight; negative is a short \\
\texttt{confidence} & number & optional ($0.5$) & stated conviction in $[0,1]$ \\
\texttt{rationale} & string & optional (\texttt{""}) & per-order rationale \\
\addlinespace
\multicolumn{4}{@{}l}{\emph{Object} \texttt{DecisionCost}} \\
\texttt{cost\_usd} & number & optional ($0$) & dollar cost of the compute spent \\
\texttt{tokens\_in} & integer & optional ($0$) & prompt tokens consumed \\
\texttt{tokens\_out} & integer & optional ($0$) & completion tokens produced \\
\texttt{reasoning\_tokens} & integer & optional ($0$) & reasoning tokens, not re-added to the total \\
\bottomrule
\end{tabular}
\end{table}

\subsection{Determinism across runtimes}\label{sec:determinism}

A scenario is a pure function of one 64-bit seed. The generator's price transform uses only multiply, add, divide and max; it never calls \texttt{ln} or \texttt{exp}, which differ across libm implementations, so a generated panel is byte-identical on any platform and any runtime. The source commits FNV-1a fingerprints of canonical generated scenarios and of a scripted order sequence's fill tape. Each fingerprint is preceded by an exact canonical-JSON fixture comparison, so a failure presents readable bytes before an opaque hash mismatch. Native, Python, host-compiled WebAssembly, \texttt{wasm32} and the committed npm binary read the relevant fixtures. Tests recompute the hashes on every commit, and WebAssembly parity tests require native-identical returns, traces and scenarios.

A separate \texttt{SPEC\_HASH} binds the eight version and source files that determine tape semantics, with CRLF normalized before an FNV-1a fold and a manual epoch for semantics outside that file set. The Rust and WebAssembly engines expose the value; the Python and npm wrappers pin it and refuse to load an absent or mismatched engine with a named diagnosis. This handshake turns a wrapper from one semantics revision paired with an engine from another into a startup error. It is a compatibility fingerprint, not a cryptographic commitment. The generator remains deliberately public and deterministic; that makes the hashes checkable and future bars computable from a recovered seed. \Cref{sec:predictability} measures the boundary: a causal AR adversary finds no undisclosed stressed-tier edge, but a public bounded band of $2^{16}$ seeds is exhaustively inverted from one observed bar in about one second. Held-out evaluation therefore requires high-entropy unpublished seeds, not merely disjoint public bands.

\subsection{Recompute-from-decisions tamper evidence}\label{sec:replay}

A rollout trace is a versioned JSONL artifact whose decision sequence is sufficient to recompute performance against the frozen dataset, costs, window and seed. Because the engine is deterministic, replaying those inputs reproduces the run byte-for-byte. The npm package's test suite exercises the adversarial case directly: it captures a trajectory, doctors every order's target weight, replays both, and asserts that the tampered decisions cannot reproduce the honest returns. The low-level replay iterator consumes decisions in order; it does not authenticate every auxiliary trace field, including stored step and observation identifiers, and it treats an exhausted decision iterator as holds. Replay therefore verifies claimed performance from decisions and frozen inputs, not the integrity or completeness of arbitrary metadata. Higher-level field and promotion validators impose their own structural checks.

\subsection{Requested versus effective configuration}\label{sec:effective-config}

A reproducible label is insufficient if the requested arm never reached the component that consumed it. The current native binding therefore reports the panel dimensions, window, execution seed and FNV-1a fingerprint of the dataset actually installed in an environment. The Python checker reads that block back, independently regenerates the panel requested by the producer and refuses any disagreement. The current F1, F3, F4, F7, predictability, sealed-seed, throughput and witness producers merge those per-seed readbacks into an \texttt{effective\_config} block when they are run. Most of the committed evidence predates this field and was not regenerated merely to add it; the check protects future regeneration and does not retroactively attest the configuration of an old artifact. F7 is the one finding regenerated since, so its artifact carries the block and the others do not.

\subsection{Procedural scenarios and disjoint seed bands}\label{sec:scenarios}

Scenario families follow Procgen's integer-seed-interval model \citep{cobbe2020procgen}: Calm, Hard and Extreme volatility-and-jump tiers, a cointegrated-pairs family with a mean-reverting spread, and a regime-shift family that hands a trend regime over to a whipsaw. The Rust core exposes \texttt{train\_test\_split}, which carves a disjoint test family (two separated integer intervals, with the disjointness asserted in tests) from a bounded train interval and refuses an unbounded one, and \texttt{cross\_regime\_split}, which holds the seed band fixed while swapping the regime, a different and generally harder probe than a within-tier split. The evaluation protocol fixes the canonical bands: train seeds $[0,256)$, an unsampled gap of 10{,}000 seeds, and a held-out test band $[10256,10512)$. The fixed-policy F3 and witness producers instead document their own 16-seed subsets: $[0,16)$ and $[10016,10032)$; the latter is a separated gap-band evaluation subset, not the canonical held-out namespace.

The Python layer additionally reserves an absolute evaluation namespace starting at seed $10^6$ and commits a named held-out regression set inside it, with a disjointness guard asserting the set never touches the train band. All of these bands are public and therefore enumerable; for an adversarial evaluation the recommended mode is the sealed derivation of \cref{sec:sealed-seeds}, which keeps the band structure checkable while hiding the concrete seeds behind a salt.

Overfitting is then a measurement. \texttt{generalization\_gap} scores the same policy on disjoint named bands and reports the differenced deflated Sharpe, and \texttt{cross\_regime\_transfer} reports the in-distribution minus zero-shot out-of-distribution gap over identical seeds, which is zero by construction when the regimes coincide.

\subsection{The task suite and the closed-form reference policy}\label{sec:tasks}

Beyond the canonical position-trading environment, the suite ships a simplex-allocation portfolio task, a VWAP/TWAP implementation-shortfall execution task, and two market models: an endogenous shared-book market where aggregate agent flow moves the cleared price through a Kyle-style linear permanent-impact term \citep{kyle1985continuous} and an Almgren--Chriss temporary-impact term \citep{almgren2001optimal}, and a deterministic limit-order-book market with integer ticks, price-time priority, a single-price call-auction uncross and a read-only walk-the-book sweep-cost query.

The two market models clear differently, and the distinction bounds what the environment can be read to say. Neither is a uniform-price mechanism. In the endogenous shared-book market every agent decides against one common reference mid for the bar, the exogenous mid scaled by a permanent-impact multiplier that accumulates only strictly earlier flow, so no order moves the price its own submitter saw; the coefficient of the Kyle-style linear permanent-impact term drives that multiplier forward. The term is a stylized linear permanent impact in the sense of \citet{huberman2004price}, not the equilibrium coefficient of \citet{kyle1985continuous}. Each agent's execution price then differs from the others', because the Almgren--Chriss temporary term charges it for both the bar's net flow and its own size \citep{almgren2001optimal}. The limit-order-book market is different again. It matches each submitted order against the resting book under price-time priority, walking as many levels as the size requires, so one step can execute at several prices and a resting order's queue position is economically real: reducing an order's size keeps its place in the level's queue and increasing it re-queues at the back. That book's single-price call-auction uncross is a separate read-only query, not the path ordinary orders take. What the two models share is discrete time. Orders arrive once per step in a canonical submission order and nothing trades between steps, which puts both closer to a frequent batch auction \citep{budish2015high} than to a continuous market and leaves latency races and the sniping of stale quotes with no representation in either. No maker-taker fee or rebate is modeled: every Sharpe, regret and markout in \cref{sec:validation,sec:findings} is gross of fees. The restriction is deliberate, it keeps clearing deterministic and byte-reproducible, but it is a scope restriction and results should be read inside it.

In the book, the canonical submission order is the submitting agent's index, so when two agents rest at the same price in one step the lower index always queues first, and seat assignment becomes a hidden treatment in any multi-agent comparison on the book. The PettingZoo order-book environment keeps that order by default and offers an opt-in rule, \texttt{priority="seeded\_shuffle"}, that draws a uniform permutation of the seats for each step from the seed and the step. Every pair of seats then queues each way with probability one half. The rule changes only the codes the Python environment submits, so the native engine, \texttt{SPEC\_HASH} and the golden fill tape are unchanged. Other simulators handle the same ordering explicitly: the market-making study of \citet[PDF p.~52]{fernandezvicente2025market} has investors pick at random among identical quotes to avoid a bias from the order in which quotes were placed, and ABIDES-MARL executes simultaneous actions in an explicitly configured order because independent agent interruptions in ABIDES-Gym cannot guarantee a consistent one \citep[p.~4]{cheridito2025abidesmarl}. A survey of language models in trading lists queue priority among the microstructure details that agent frameworks leave unstandardized \citep[p.~12]{fu2025newquant}.

The same environment values each agent's inventory at a mark that excludes the agent's own orders. By default (\texttt{mark="ex\_own\_mid"}) the mark is the midpoint of the best bid and best ask among resting orders that other controllers own, so an agent's own resting quotes never enter its own valuation; the earlier whole-book mid remains available as \texttt{mark="book\_mid"} to replay earlier runs. Each seat is its own controller unless the \texttt{controllers} argument names the seats one entrant runs, and seats that share a controller neither value each other's quotes nor set each other's mark by trading together. The principle matches CoffeeBench, which values consumed inventory at a weighted-average cost basis so that an agent cannot inflate its score by self-trading at marked-up prices \citep[p.~5]{sugiura2026coffeebench}. The mark has limits. An agent's quotes still move the reference mid that every seat quotes around, so its rivals' next quotes, and with them its own mark, can follow its quotes. While the other controllers lack a resting order on either side, the mark is the price of the step's last fill between two different controllers, and it carries its last value only on a step with no such fill. A lone agent's only counterparty is the noise trader, so it is marked where it last traded and its reward moves only on the steps where it trades. Over seeds 0 to 31 with no inventory penalty, a lone agent that posted its bid one tick below the reference mid and its ask twenty ticks above it, walking the mid for sixty steps before quoting one tick on each side, earns a mean episode reward of 881.4 against 897.1 for its cash plus inventory at the final reference mid, where the frozen mark it replaced paid 2360.3. Under the default penalty, and with the skewed quote held for the whole episode, the mark ranks that agent below a symmetric three-tick quoter over the same seeds, as \texttt{mark="book\_mid"} does. The frozen mark ranked it above. The book also has no self-trade prevention; a self-trade leaves one agent's cash and inventory unchanged.

The endogenous market's impact coefficients $(\lambda,\eta)$ are point estimates, and an opt-in elliptic uncertainty set replaces the pair with an ellipse around it, intersected with the nonnegative quadrant, whose two half-widths and correlation are declared modelling inputs rather than estimates the engine validates. The robust entry points clear each bar at the most expensive coefficients inside that set, found in closed form, and given no set they execute the point-estimate arithmetic byte for byte, which the native equality test named in \cref{sec:findings} pins. The construction is adapted from \citet{ma2025robust}, whose elliptic sets are defined around several foci in the space of transition-kernel perturbations (p.~5) rather than over two impact coefficients, and who report that symmetric uncertainty sets often produce overly conservative strategies (p.~9). A rank-neutral report, \texttt{impact\_misspecification\_gap}, runs one policy alone on each seed twice on the linear kernel, once at the point estimate and once against a declared set. It refuses with a typed error unless both arms traded on the same exogenous path, and it returns per-seed return and per-bar Sharpe under each arm with t-based intervals on their differences. The difference adapts Ma and Huang's robustness measure, the relative portfolio gap with and without market impact (p.~9), to a point-estimate-versus-worst-case pair, and it does not reduce to a cost of misspecification. Both arms mark the held position at the exogenous mid, so neither return carries its own permanent impact, and each row reports what the engine's cleared-mid valuation would add as a separate own-impact mark. With a positive $\lambda$ radius the worst-case coefficient still moves the cleared mids, so the two arms' weights, order sizes and later fills can differ, and the gap has no guaranteed sign. Only an $\eta$-only set guarantees a nonnegative gap, and only for a policy that sends identical weights in both arms while cleared mids stay positive. The two premises are checked rather than assumed. A cleared mid that is not positive refuses the arm with a typed error, and every seed carries \texttt{identical\_actions}. The report's \texttt{sign\_guaranteed} is true only for an $\eta$-only set whose policy sent identical weights on every seed. No result in this paper uses the uncertainty set or the report.

The market-making task is the suite's calibration instrument. It implements the Avellaneda--Stoikov model \citep{avellaneda2008high} and ships the model's closed-form quoting policy beside the environment: reservation price $r = S - q\gamma\sigma^2\tau$ around the mid price $S$ and half-spread $\delta^* = \gamma\sigma^2\tau/2 + \gamma^{-1}\ln(1+\gamma/\kappa)$, skewed by inventory. This closed form is the source model's asymptotic approximation, not an exact optimum: the exact treatment of the same problem is given by \citet{gueant2013dealing}, and no proof is offered that the policy is optimal for the implemented reward, which adds an inventory cap, a running inventory penalty and a terminal liquidation charge the source model lacks. We therefore call it the closed-form reference policy. Both the environment and the policy are pure functions of the same frozen parameter set, so the regret metric \texttt{mm\_regret} runs the candidate and the reference on identical seeds and reports the mean reward gap; the reference scores zero against itself by construction, which fixes the metric's zero point without claiming the zero point is unbeatable. Identical seeds pair the two arms only while both consume the environment's random stream identically. Arrivals, fills and the mid step come from one generator per episode, and numpy's binomial sampler takes a number of draws that depends on the fill probability, so two quoting policies can desynchronize the stream, after which their arrivals and mid paths differ and the gap mixes policy with luck. \texttt{mm\_regret} therefore records the mid, the two arrival counts and the generator state after every step, and raises \texttt{UnpairedMidPathError} at the first episode where any of the three differs between the arms. Mids alone would not show the desynchronization, because at zero volatility, or at a volatility whose step falls below the float resolution of the mid, every mid equals the opening price while the arrivals differ. The default stream is unchanged, and a committed test reproduces every regret of \cref{sec:f2} bit for bit through the checked function, so each of those comparisons drew the same stream. The check does not couple the two arms' fill draws, which depend on each arm's own quotes. ABIDES avoids the problem by giving each agent its own generator seeded from one seed, so that an A/B comparison leaves the other agents' random streams unchanged (as described by \citealp[PDF p.~41]{fernandezvicente2025market}).

Each step also splits its reward into four parts that sum to it up to rounding: spread capture on the step's fills, inventory marked to the mid move, the running inventory penalty, and the terminal liquidation cost. \texttt{mm\_pnl\_split} averages that split per policy over the episodes \texttt{mm\_regret} uses, which separates a quoter that earns spread from one that earns by holding inventory into a favourable path. The split adapts the four-term reward of \citet[Eq.~4.1, PDF p.~67]{fernandezvicente2025market}; this environment has no hedging cost, and its running penalty and liquidation charge take the place of that reward's inventory penalty. Differencing two policies' splits attributes their regret component by component only when both arms drew the same stream, which \texttt{mm\_pnl\_split} does not itself check; \texttt{mm\_regret\_split} runs both arms through the check above and returns the per-component regret. Neither output feeds a score. The model also assumes fills that do not select against the maker, which sits in deliberate tension with the adverse-selection probe of \cref{sec:f6}. Gym-style market-making environments built on the same model family exist \citep{jerome2023mbtgym}; what this task adds is the shipped reference policy and the regret-not-PnL scoring convention inside the verifiable-producer stack.

\subsection{Mandates and deferred claims}\label{sec:mandates}

Each scenario can sample a trading mandate, and the objective is conditioned on it. A mandate pairs a structural style with up to three optional constraints layered on top of that style: a realized-drawdown cap, a per-bar gross-exposure cap on $\sum_i |w_i|$, and a benchmark symbol to beat. The wire vocabulary names five styles, long-only, market-neutral, momentum, unconstrained and pairs-convergence, and four of them are drawn. Unconstrained is the permissive control, and the two styles that need a short leg, market-neutral and pairs-convergence, are excluded when the environment disallows shorting, so a sampled mandate is never unsatisfiable by construction: \texttt{mandate\_breach} detects structural violations from the realized weights and returns, and the training rubric penalizes wrong-objective behavior rather than merely not rewarding it. Momentum is the style that is not drawn. It renders ``lean into recent winners, cut losers'' in its prompt text, and the breach checker sees per-bar portfolio weights and pooled returns and never per-symbol returns, so nothing in the grader can read that constraint off a book. An objective no grader can read must not be the objective an episode is scored against, so the label stays in the vocabulary, because recorded traces carry it and must still parse, and is absent from the draw. A deterministic failure taxonomy classifies every episode into a single worst-case-first disposition, cascade-wiped before bankrupt before stopped-out before mandate breach before clean, and a suite-level rollup tallies the distribution with a schema that always carries every mode. Two further dispositions describe the evidence rather than the behavior: an episode whose returns, events or mandate are malformed or missing is invalid evidence, and one whose event stream records a protocol, transport, agent or infrastructure error is an execution failure. Neither can count as clean, and neither attributes an infrastructure fault to the agent. The mandate checker, the reward schemes, the reward-misspecification controls and the taxonomy read the per-bar weight vectors through one shared event contract, which refuses a weight payload sent under any event name but the canonical one rather than letting the readers disagree about it. A cell that the checked Rust entry point refuses at the wire (\cref{sec:transport-gate}) is different again: it produces no episode to classify, carries a typed transport fault instead of a disposition, and is counted apart from the rollup that \cref{sec:f7} reports.

Questions that resolve after the episode get the same structural treatment as look-ahead. A \texttt{DeferredDesk} is the only object the agent touches: its entire state is a list of claims and one integer clock advanced by the harness, it holds no dataset and no environment reference, and it computes no score, so there is no code path from a claim to its resolving datum. Resolution is a free function over claims and outcomes that refuses any outcome available at or before the claim's commit bar and any outcome predating the stated horizon. The agent cannot see the answer at commit time because the object it holds cannot produce one.

\subsection{The probe layer: distrusting the simulator}\label{sec:probes}

\textbf{Realism.} A stylized-facts battery computes excess kurtosis, absolute-return autocorrelation, gain/loss skew and aggregational Gaussianity in the sense of the classical catalogue \citep{cont2001empirical}, plus a timescale-asymmetry statistic in the sense of \citet{zumbach2009time} and an exploratory Fano exceedance-count ratio (the variance-to-mean ratio of the number of large-return exceedances per equal-length block, which is 1 for an IID Poisson-like count and above 1 when exceedances bunch). \texttt{certify\_realism} gates kurtosis, an ACF value above the fixed 95th-percentile IID null for its exact short-panel estimator, and movement of absolute excess kurtosis toward zero on aggregation. Fano is reported but ungated: it is normalized by its exact conditional-IID finite-panel expectation, yet 12 count windows remain too noisy for a certificate. Certification is reported beside the results rather than required for the leaderboard, and the canonical configuration fails it on 23 of 24 seeded panels (\cref{sec:f4}). The intended endgame is a generator revision that passes on disjoint calibration seeds or leaderboard claims conditioned on certification; until then the gate keeps the failure visible.

\textbf{Manipulation.} A crude push-and-unwind schedule, the classic trade-based manipulation of \citet{allen1992stock}, is scored against its own zero-impact reference: the live and reference runs share the seed and follower population, and differ only in impact. A boundary sweep varies one parameter and a size-response sweep asks whether payoff is bounded. Linear permanent impact makes the baseline a model-consistency check under the assumptions of \citet{huberman2004price,gatheral2010no}. On separate entry points, nonlinear permanent impact is $\lambda\operatorname{sign}(Q/V)|Q/V|^{\beta}$, with dimensionless flow $Q/V$ and a dimensionless coefficient $\lambda$ for each declared $\beta$; $\beta=1$ uses the exact original linear arithmetic. This avoids overloading the market-making model's $\kappa$ (order-book-depth parameter). \Cref{sec:f5-concave} evaluates the symmetric schedule, and \cref{sec:f5-positive-control} reports a 135-cell exploratory asymmetric search with selection-adjusted inference. Each probe result now also carries a per-follower seat-removal externality: a follower's P\&L in the live market with the manipulator trading its schedule, minus the same follower's P\&L on the same seed with the manipulator held flat for the whole episode. Every entry is exactly zero when the push weight or the permanent-impact coefficient is zero. The device adapts the competency measurement of EconEvals, which removes the competitor to turn a two-seller pricing game into a single-agent one \citep[p.~61]{fish2026econevals}; here the manipulator is held flat to measure what its presence does to the other seats. The externality is reported beside the attribution, feeds no rank, gate or boundary, and postdates the committed F5 artifact.

\textbf{Adverse selection.} Informed meta-order flow is routed against a maker roster, with fills struck at the pre-move mid so the maker's markout carries the informed trader's edge. The mechanism the probe operationalizes is the classical one \citep{glosten1985bid}, with one deliberate simplification: the price path is exogenous and quotes do not widen in equilibrium response to toxicity, so the probe measures markout accounting against drift, not equilibrium adverse selection. The design is a paired control: the informed leg sides child orders on the alpha, the uninformed leg sides the identical flow on a coin while holding the price path fixed, and the comparison reports mean markout per filled unit at each horizon. If the informed leg does not mark out worse, the module's own numbers are declared noise.

\textbf{Ecology.} A replicator dynamic in the market-ecology tradition \citep{farmer2002market,lux1999scaling} runs a population of strategies over generations: fitness earned in a shared-book market where a species that gains share becomes a larger fraction of the flow everyone else fills against, seat allocation from population shares, extinction below a threshold, an innovation operator that breeds parameter-perturbed variants, and shock schedules that rotate regimes or scale impact to simulate liquidity droughts. The control schedule holds the environment steady, so any claim about a shock has a baseline to beat.
